\documentclass[11pt]{article}

\usepackage[T1]{fontenc}
\usepackage[utf8]{inputenc}
\usepackage[a4paper,margin=23mm]{geometry}
\usepackage{booktabs}
\usepackage{longtable}
\usepackage{array}
\usepackage{tabularx}
\usepackage{xcolor}
\usepackage{listings}
\usepackage{hyperref}

\hypersetup{
  colorlinks=true,
  linkcolor=blue!55!black,
  urlcolor=blue!55!black,
  pdftitle={Reducing Memory in Hint-Guided Prover9 Saturation},
  pdfauthor={Prover9 engineering report}
}

\lstset{
  basicstyle=\ttfamily\small,
  breaklines=true,
  columns=fullflexible,
  frame=single,
  framesep=5pt,
  rulecolor=\color{black!25},
  showstringspaces=false
}

\newcommand{\setting}[1]{\texttt{\detokenize{#1}}}
\newcommand{\hashval}[1]{\nolinkurl{#1}}
\newcommand{\mib}{\,MiB}
\newcommand{\kib}{\,KiB}
\newcolumntype{Y}{>{\raggedright\arraybackslash}X}

\title{Reducing Memory in Hint-Guided Prover9 Saturation\\
\large From Passive-Clause Explosion to a Compact
OTTER-Compatible Given-Clause Loop}
\author{Technical engineering report}
\date{11 August 2026}

\begin{document}
\maketitle

\begin{abstract}
Long hint-guided Prover9 searches in algebra can retain millions of clauses in
the set of support (SOS), millions of demodulators, millions of disabled proof
ancestors, and hundreds of thousands of hints.  Historical runs reported
18.8\,GiB for an Osborn proof and 125.8\,GiB for an AAPERM prefix.  The initial
hypothesis that disabled clauses alone caused this growth was only partly
correct.  Clauses rejected immediately as tautologies or by forward
subsumption were not retained, while deleting almost all disabled clauses saved
only 22\% at the 4000-given Osborn boundary.  The dominant architectural
problem was that retained but not yet selected SOS clauses remained represented
as ordinary pointer-rich clause and term graphs and participated in eager
simplification indexes.

This report describes two responses.  The first is a DISCOUNT-style loop with a
strict active/passive separation, dense passive records, packed hints, and a
bounded collective representation of delayed paramodulation and
hyperresolution work inspired by Waldmeister.  It gives the strongest
asymptotic memory bound, but changes when simplification and generating
inferences occur.  On the supplied Osborn hint chains, matching more hints did
not reproduce the historical proof: demodulation and inference timing are part
of the effective search strategy.

The accepted response is an eager compact OTTER-compatible loop.  It preserves
the old treatment of every retained clause for demodulation, unit conflict,
subsumption, backward demodulation, hint matching, and given-clause selection,
but replaces resident clause graphs and pointer-based passive indexes by an
immutable file archive, dense selector records, stable clause numbers, shared
serialized terms, and four compact authoritative indexes.  Ordinary clauses
are reconstructed only for exact logical tests or activation.

On the accepted full proof, old P9 used 550,400\kib peak RSS and 765.69 seconds
of user CPU.  Compact OTTER proved at the same 2945-given boundary, with the
same 7051 proof clauses and 3231 new hints, in 754.78 user seconds and
127,420\kib peak RSS.  This is a measured 76.85\% whole-process reduction
(4.32 times smaller) with 1.42\% less user CPU.  A terminal accounting explains
96.93\% of proportional resident memory.
\end{abstract}

\paragraph{Implementation status.}
The implementation is on branch \setting{phase5-compact-frontier}.  The
accepted code commit is \setting{2483a7a49e4e21c54b1692c142b55435eeefe44b};
the closing audit is commit \setting{ac3f0fe}.  All hard Phase-5 correctness,
proof, CPU, memory, checkpoint, and accounting gates passed on the current
\setting{chat_test}/Osborn acceptance problem.

\section{Terminology and scope}

The report uses the standard saturation terminology familiar from OTTER,
Prover9, E, Vampire, the \emph{Handbook of Automated Reasoning}, and Harrison's
\emph{Handbook of Practical Logic and Automated Reasoning}.

The \emph{passive set}, called \setting{sos} by Prover9, contains retained
clauses waiting to be selected.  The \emph{active set}, called
\setting{usable}, contains selected clauses available as parents of generating
inferences.  A \emph{given clause} is selected from passive, moved to active,
and combined with suitable active clauses by the enabled inference rules.
Generating inferences are called expansion; simplification and redundancy
elimination---demodulation, unit deletion, forward/backward subsumption, and
backward demodulation---are called contraction.

A hint is a clause used for search guidance, not an inference premise.  A
derived clause that subsumes a hint, is subsumed by a hint, or satisfies the
configured matching relation can receive favorable selection treatment.  This
is the hints/proof-sketch tradition developed by Veroff and implemented in
OTTER and Prover9.

The main nonstandard terms are representational:
\begin{itemize}
  \item A \emph{stable clause number} is the ordinary Prover9 clause ID used in
        an index instead of a C pointer.
  \item A \emph{dense record} is one fixed-size array element holding selection
        and status data rather than a separately allocated object graph.
  \item \emph{Materialization} reconstructs an ordinary Prover9 clause,
        literal, and term graph from an immutable serialized record for an
        exact test.
  \item A \emph{posting list} is a list of stable clause numbers sharing an
        index feature such as a root symbol or shallow term path.
  \item A \emph{conservative filter} may return false candidates but never omit
        a possible answer.  The established matcher, unifier, subsumption test,
        or rewrite operation makes the final decision.
\end{itemize}

Accordingly, ``exact index'' below means that the final operation agrees with
the ordinary implementation.  The intermediate candidate set need not itself
be exact.

\section{The original memory problem}

The motivating AIM searches are much larger than their given counts suggest.
A few thousand selected clauses can generate and retain millions of SOS
clauses, and many of those retained equations can become simplifiers before
selection.

\begin{center}
\small
\begin{tabular}{lrrrrr}
\toprule
State & Given & SOS & Demodulators & Disabled & Hints \\
\midrule
Osborn, 4000 prefix & 4,000 & 588,252 & 489,224 & 577,347 & 310,153 \\
Osborn proof & 11,242 & 4,594,786 & 3,759,850 & 6,519,452 & 310,153 \\
AAPERM prefix & 12,560 & 8,504,637 & 6,199,295 & 4,593,127 & 220,117 \\
\bottomrule
\end{tabular}
\end{center}

The corresponding internal \setting{Megabytes} values were 2626.98, 18810.85,
and 125829.01\mib.  The AAPERM profile further showed that inference generation
was not the dominant CPU term: about 902 of 225,533 user seconds were in
\setting{infer}; preprocessing, subsumption, indexing, disabling, and backward
demodulation dominated.

\subsection{What the 2024 RAM discussion corrected}

The archived \setting{Plans-RAM-24.txt} discussion prevents two incorrect
interpretations.

First, \setting{Generated} is not a resident-clause count.  Immediate
tautologies and clauses rejected by forward subsumption are normally destroyed.
They consume CPU and allocation traffic but not terminal memory by their count.

Second, disabled clauses matter, especially late in a long run, but are not the
whole problem.  Direct deletion experiments measured:

\begin{center}
\begin{tabular}{lrrr}
\toprule
Boundary & Ordinary & Delete disabled & Saving \\
\midrule
1000 givens, hints & 674\mib & 644\mib & 4.5\% \\
2000 givens, hints & 1101\mib & 963\mib & 12.5\% \\
4000 givens, hints & 2627\mib & 2037\mib & 22\% \\
8000 givens & 9192\mib & 6470\mib & 30\% \\
\bottomrule
\end{tabular}
\end{center}

At 4000 givens the disabled population fell from 577,347 to 383 while the
588,252-clause SOS and its eager indexes remained.  Moreover, proof
justifications refer to parent clause numbers.  Freeing a disabled object
without a replacement proof-history representation leaves dangling
references.  A radical solution therefore had to compact proof history and the
live SOS/index state.

\subsection{Logical bytes are not resident bytes}

Prover9's historical \setting{Megabytes} is allocator-oriented and may omit
anonymous mappings, mapped files, libc arenas, shared pages, and file-backed
resident pages.  Conversely, the logical length of an archive file is disk
usage rather than RAM.  Measurements in this work distinguish internal
component bytes, RSS, PSS, and the maximum RSS from \setting{/usr/bin/time -v}.

This mattered in a long DISCOUNT run whose deleted mmap file was
13,419,485,217 logical bytes, occupied 9.7\,GiB of allocated disk blocks, and
had 10,223,128\kib actually resident.  An mmap does not stay cold if code scans
it.  A former statistics routine verified every archived record and thereby
faulted the entire mapping into memory.  The accepted configuration uses
explicit file reads instead.

Explicit I/O can still populate the kernel page cache.  Those pages are
reclaimable and absent from process RSS/PSS, but may be charged to a cgroup and
consume machine RAM until reclaimed.  Thus 127,420\kib is an exact process
peak-RSS result, not every kernel page used on behalf of the job.  Total
container or machine claims require cgroup \setting{memory.current},
\setting{memory.peak}, and the \setting{file} fields of
\setting{memory.stat}.

\section{Semantic constraints}

A Prover9 SOS clause is passive for generating inferences but, in the
traditional OTTER loop, not passive for contraction.  A newly retained equation
can become a demodulator immediately, rewrite old SOS and usable clauses, rewrite
hints, and change later subsumption decisions.  Unit clauses participate in
unit conflict and simplification; nonunit clauses must remain visible to
subsumption.  Hint matching occurs while a candidate is an ordinary clause,
before passivization, and affects its weight and selector membership.

Four contracts follow.

\paragraph{Search.}
Trajectory-preserving mode must execute the same eager contraction transaction
for a retained clause.  Selection cycles, weights, hint degradation,
\setting{match_once}, matcher limits, and actions must observe the same result
in the same order.

\paragraph{Hints.}
Hints are not active inference premises, but equivalence handling, ordinary and
flipped matching, subsumption matching, degradation, matcher limits,
\setting{back_demod_hints}, and checkpoint state must remain exact.  A cold
passive is checked against hints while it is still a materialized candidate and
is archived only after authoritative hint processing.

\paragraph{Proofs.}
Every parent must remain recoverable by stable number after its ordinary body
is freed.  The archive retains clause body, attributes, flags, matching-hint
state, and justification.  Corruption must be detected rather than appear in a
proof.

\paragraph{Indexes.}
Removing passive clause graphs also removes pointers stored in ordinary
indexes.  Every passive-sensitive operation therefore needs either a compact
authoritative replacement or an ordinary live representation.  Compact OTTER
refuses to start unless all four required replacements are active.

\section{Waldmeister and the DISCOUNT/collective path}

Hillenbrand's Waldmeister account emphasizes compact discrimination structures,
a strict active/passive separation, and collective representation of critical
pairs.  Critical pairs generated with one newly active equation can be stored
as a constant-size descriptor plus a minimum remaining weight, while only a
fixed buffer of promising pairs is represented individually.  This changes
Waldmeister's passive-space growth from quadratic to linear in abstract search
time.  Its reported unit-equational workload processed more than 500 million
passive equations with about 200 MB.

That number cannot be transferred directly to Prover9's paramodulation and
hyperresolution workloads.  The architectural question is nevertheless the
right one: represent delayed inference work rather than every conclusion.

\subsection{Strict active/passive ownership}

\setting{assign(search_loop,discount)} inserts only selected active clauses in
inference and simplification indexes.  A new candidate is simplified,
hint-matched, weighed, and stored passively.  At selection it is refreshed
against the current simplifier epoch, requeued if changed, and only then becomes
an inference parent or demodulator.

\setting{assign(passive_store,dense)} replaces each passive clause graph and
selector nodes by a dense record containing stable ID, selector keys, weight,
semantics, matched-hint ID, and epoch/status bits.  The clause and justification
are serialized and reconstructed at selection.

\subsection{Packed hints}

The old FPA hint index retained pointer-rich term forests.  The packed modes
retain compressed hint bodies, structural postings, and small mutable side
arrays.  The first packed implementation saved RAM but was too slow: at one
1000-given boundary it used 295,680\kib versus old P9's 507,108\kib, but took
441.39 instead of 113.44 CPU seconds.  Instrumentation exposed broad posting
lists and repeated reconstruction.

The successor \setting{packed_fast} uses density-adaptive intersections, exact
feature profiles, a bounded dependency-scoped cache, direct compressed-unit
matching, and compact \setting{_AnyConst} state.  At a 300-given selected-
DISCOUNT boundary it took 10.23 user seconds versus 27.72 for FPA and 50.02 for
the earlier packed mode, with a byte-identical hint trace.

\subsection{Collective expansion}

The collective frontier records bounded descriptors for paramodulation from
and into a given clause and for positive and negative hyperresolution.  Each
descriptor holds stable parent IDs, an active-set snapshot boundary, a native
rule continuation, and fairness/checkpoint state.  A balanced scheduler gives
each enabled rule a nonzero share, imposes high/low descriptor water marks, and
bounds raw conclusions committed per turn.

Optional read-only hint previews rank a bounded candidate set without assigning
an ID or mutating hints.  A dual-cursor discovery lane can promote a promising
conclusion early and later verify by ordinal and structural checksum which fair
conclusion to skip.  FIFO service remains, supplying the fairness argument.

\subsection{Why it did not reproduce Osborn}

Some DISCOUNT/collective runs matched more hints than the historical proof but
did not prove within the tested resources.  A hint match says that a useful
clause shape appeared; it does not say that the same ancestors exist, the same
demodulators were available, or the clause survived in the same normal form.
Traditional OTTER admits an unselected SOS equation as a demodulator
immediately.  Selected DISCOUNT delays it.  Eager-interreduced DISCOUNT rewrites
more strongly but can remove shapes expected by the historical chain.

Bounded rewrite-debt drains fixed a liveness defect in which eager repair could
exclude all ordinary inference, but did not recover the old trajectory.  The
result illustrates a familiar ATP fact: fairness or refutational completeness
does not imply the same proof under a finite resource limit.  Demodulation and
inference scheduling are part of the practical strategy.

\section{The accepted eager compact-OTTER architecture}

Phase 5 retains \setting{search_loop=otter} and
\setting{inference_frontier=clauses}.  Each retained candidate completes the
ordinary eager transaction.  Afterwards its body is archived and long-lived
passive structures retain compact records by stable number.

\begin{lstlisting}
infer an ordinary candidate
  -> demodulate and simplify literals
  -> safe unit-conflict test
  -> exact hint matching and weighting
  -> limits, semantics, keep/delete rules
  -> forward subsumption
  -> assign a stable clause number
  -> admit a demodulator when appropriate
  -> backward subsumption and backward demodulation
  -> update the four compact indexes
  -> serialize body and justification
  -> retain only dense selection state and stable IDs

select a given clause
  -> locate its archive record by stable number
  -> reconstruct the ordinary clause
  -> retire compact passive memberships
  -> activate the clause
  -> make ordinary inferences at the historical point
\end{lstlisting}

\subsection{Immutable file archive}

\setting{ancestor_store=file} uses checksummed append-only records and explicit
\setting{pread}/\setting{pwrite}.  A record retains the compressed clause,
attributes, flags, hint state, and parent IDs.  In compact OTTER the proof
archive is also the sole owner of a cold passive body; the selector does not
retain a duplicate arena.

The accepted proof's archive was 84,404,096 logical bytes and used one 4-KiB I/O
buffer.  It performed about 3.8 million reads totaling 374 MB and 412,000 writes
totaling 84.4 MB.  Logical file length is reported separately from process RAM;
reclaimable kernel page cache needs cgroup accounting for a whole-job measure.

\subsection{Dense given-clause selection}

Dense records preserve low/high selection cycles, breadth-first level,
semantics, weight, hint ID, age, and membership.  At the final boundary 131,001
passives used 9,532,864 bytes of records and a 611,368-byte selector heap.

A set-associative decoded-passive cache is available, but the accepted setting
is zero.  Controlled 0, 4, and 32-MiB experiments at 1500 givens avoided about
98,000 decodes without improving CPU; nonzero caches increased RSS.  Index
candidate work, not archive decoding, was the dominant cost.

\subsection{Four authoritative compact indexes}

The compact rewrite bank stores serialized rule sides, rule kind, activity, and
stable proof ID.  A radix-compressed discrimination structure retrieves rule
candidates; reverse occurrence streams support backward demodulation.  Exact
matching, ordering, and rewriting retain ordinary Prover9 semantics.  It used
11,205,308 bytes at the proof boundary.

The compact unit index supports generalization, instance and unifier retrieval,
forward/back unit subsumption, and unit conflict.  It stores radix paths,
postings, polarity, a symbol mask, and stable proof ID, using 13,166,960 bytes.

The backward-demodulation index groups occurrences by root symbol, clause
number, and exact shallow path signature.  Monotone record and occurrence
positions are delta-encoded in small blocks.  A structural matcher checks every
survivor.  It used 12,109,376 bytes.

The nonunit feature trie returns stable IDs for forward/backward subsumption;
ordinary subsumption remains authoritative.  It used 1,508,576 bytes.

\subsection{Shared serialized terms}

The four indexes share a clause-coalescing 32-bit term-token pool.  A clause is
serialized once; subsequent index requests receive stable slices.  The pool has
a sparse stable-ID directory, 12.5\% growth, and coordinated compaction.  Exact
old-to-new offset maps rebuild every dependent index before commit.  Large
rebase ordering is streamed and file-sorted, and Linux can use
\setting{mremap} to move page tables instead of copying the entire arena.

At the proof boundary the pool held 3,220,436 logical tokens in 16,219,320
bytes.  Five compactions reclaimed 14,793,444 bytes, and 38,557,344 token
requests were reused across indexes.

This is not full hash-consing.  Instrumentation at 1000 givens found 728,510
subterm occurrences but only 91,686 unique terms.  A canonical variable-record
DAG was estimated at 1,304,344 logical bytes versus 4,194,304 bytes of allocated
token capacity at that experiment.  A production hash table, fingerprints,
reference management, and slack would reduce the saving.

\subsection{Packed-fast hints}

The accepted hint bank combines compressed immutable bodies, stable-ID
postings, mutable side state, and a bounded exact query cache.  Profiles with at
most eight keys can use the 16,384-entry cache; wider profiles take the ordinary
authoritative intersection path.  Entries are 136 bytes, for a 2,228,224-byte
table.  The final hit rate was 35.08\%, avoiding 136,373,547 posting candidates.

Hint rewriting and reindexing, degradation, match-once state, and checkpoints
remain active.  Hint nodes, references, tables, cache, and bodies together used
about 15.2\mib for 88,494 hints.

\subsection{Reclamation and peak control}

Steady-state compactness did not by itself control peak RSS.  The final design
also:
\begin{itemize}
  \item rebuilds stale compact indexes at a deterministic percentage;
  \item reclaims the term pool after a measured stale-payload threshold;
  \item streams back-index rebuilds in 4096-ID batches;
  \item file-sorts large retained-ID/rebase orderings;
  \item releases predecessor arrays before allocating successors;
  \item reuses packed records during rebuilds;
  \item replaces general hashes by sparse direct maps where possible;
  \item releases dense passives, hints, and compact search indexes before final
        proof-DAG expansion;
  \item selects a compact process heap at startup with
        \setting{P9_COMPACT_HEAP=1}.
\end{itemize}

These lifetime changes prevent a compaction or terminal proof from holding two
large generations simultaneously.

\section{Correctness evidence}

\subsection{Differential and focused tests}

Focused tests cover archive corruption, allocator lifetime, dense selection,
term rebasing, stable-ID maps, rewrite retrieval, unit operations,
backward-demodulation candidates, nonunit features, hint postings/previews,
compressed-unit matching, and resumable paramodulation/hyperresolution.  The
aggregate commands are:

\begin{lstlisting}[language=bash]
make all -j4
make test1
make compact-frontier-tests
make discount-tests
\end{lstlisting}

\subsection{Exact prefixes}

At 300 givens the final file build emits 132,567 candidate, hint, kept, and
given events byte-identically to the accepted full-body \setting{packed_fast}
OTTER reference:

\begin{lstlisting}
SHA-256 bc38f369ef02271d9b8a00db0cd01a6ba8f890e9608e0abd947de68f763322aa
Given=301 Generated=120793 Kept=5737
Usable=285 SOS=4298 Demodulators=3282 Disabled=1183
\end{lstlisting}

A current 1000-given A/B comparison measured 76.90 user seconds for compact
indexes with resident bodies and 81.35 seconds for the file archive with zero
cache, a ratio of 1.058.  Both peaked at 90,404\kib and ended at
\setting{Given=1001}, \setting{Generated=1268285}, \setting{Kept=33909},
\setting{Usable=949}, \setting{SOS=26052},
\setting{Demodulators=21741}, and \setting{Disabled=6937}, with identical hint
and index counters.

\subsection{Checkpoint and proof}

A full-hint run checkpointed at given 100 and resumed to given 301.  Its
concatenated event stream has the same hash as the uninterrupted run; all
populations agree and 22 integrity checks pass.  Format 3 now records the count
of deliberately omitted, dead ID-0 scratch clauses without serializing their
bodies.

The accepted complete run ends with:

\begin{lstlisting}
Given=2945 Generated=8248032 Kept=272787 proofs=1
Usable=1800 SOS=131001 Demodulators=109987 Limbo=300 Disabled=139715
Hints=88494 Active_Hints=4122
\end{lstlisting}

\setting{prooftrans parents_only} emits 7051 proof clauses and 3231 new hints.
The normalized proof SHA-256 is
\hashval{9d7c9a12894c1c11ede6aeae08d1cec658ccee66a47fb9663859347a5413fd07},
identical to the full-body packed-fast reference.  Old FPA retains two extra
\setting{other} clauses, shifting later IDs, but reaches the same given boundary,
SOS/demodulator populations, proof length, and theorem.

\section{Evaluation}

\subsection{Accepted proof result}

\begin{center}
\begin{tabular}{lrrr}
\toprule
Measure & Old P9 OTTER/FPA & Compact OTTER & Change \\
\midrule
Given & 2,945 & 2,945 & same \\
Generated & 8,248,034 & 8,248,032 & $-2$ \\
Kept & 272,789 & 272,787 & $-2$ \\
Proof clauses & 7,051 & 7,051 & same \\
New proof hints & 3,231 & 3,231 & same \\
User CPU & 765.69 s & 754.78 s & $-1.42\%$ \\
Wall & 14:16.20 & 14:19.48 & essentially equal \\
Peak RSS & 550,400\kib & 127,420\kib & $-76.85\%$ \\
\bottomrule
\end{tabular}
\end{center}

The new process is 4.32 times smaller.  It passes the 128,000-KiB hard gate by
only 580 KiB, so a different libc, compiler, or machine requires remeasurement.
The optional 115-MiB stretch target is not met.

\subsection{Resident accounting}

\begin{center}
\small
\begin{tabular}{lr}
\toprule
Component & Bytes \\
\midrule
Rewrite bank & 11,205,308 \\
Unit index & 13,166,960 \\
Backward-demodulation index & 12,109,376 \\
Nonunit index & 1,508,576 \\
Shared term pool & 16,219,320 \\
Dense passive records and heap & 10,144,232 \\
Hint nodes, references, tables/cache, bodies & 15,953,167 \\
Ancestor handles & 4,198,608 \\
Clause-ID structures & 2,316,400 \\
Prover9 allocator reservation & 24,123,712 \\
Nonanonymous pages & 1,025,024 \\
\midrule
Accounted total & 111,970,683 \\
\bottomrule
\end{tabular}
\end{center}

The total is 106.78\mib.  Terminal PSS was 112,809\kib (110.17\mib), giving
96.93\% coverage.  Sampled RSS peaked at 123,084\kib and external maximum RSS at
127,420\kib.  The 84,404,096-byte logical file length is excluded from process
RAM; reclaimable page cache is outside this PSS accounting.

\subsection{Within-design index reduction}

At 1000 givens the first four compact indexes occupied 44,545,464 bytes.  Radix
paths, shared terms, delta streams, dense maps, packed records, and tighter
growth reduced the corresponding five-structure total to 13,239,168 bytes, a
70.28\% reduction with unchanged search state.  This explains why serializing
bodies alone was insufficient: at that boundary the early compact indexes were
about 43 MB while the clause archive was only about 8 MB.

\subsection{Newly supplied 11,000-given comparison}

The new files \setting{bob/chat_test.new.out1.gz} and
\setting{bob/chat_test.new.out2.gz} use a different, larger profile with 18,306
hints.  Both prove after about 11,000 givens.  The first is ordinary OTTER/FPA
with full bodies; the second is an early dense/packed-fast/all-compact-index
mmap configuration.  It predates the final filtering, bounded rebuild,
file-archive, cache-sizing, transient-release, and heap work.

The compressed output SHA-256 values are
\hashval{fd7a07f0a8ec1feafacb2672b664a393368ce728d0918c5e2708c7919d4a4406}
for \setting{out1} and
\hashval{049a44cb9871fb28c3cf829011fb102dc390d5101ed28677dd7ced15ac5d64f2}
for \setting{out2}.

\begin{center}
\small
\begin{tabular}{lrrr}
\toprule
Measure & Ordinary & Early compact & Change/ratio \\
\midrule
Given & 11,368 & 11,369 & $+1$ \\
Generated & 253,338,893 & 253,302,129 & $-36,764$ \\
Kept & 2,216,836 & 2,207,014 & $-9,822$ \\
SOS & 1,454,686 & 1,520,775 & $+4.5\%$ \\
Demodulators & 1,297,818 & 1,362,847 & $+5.0\%$ \\
Disabled & 739,073 & 662,309 & $-10.4\%$ \\
Internal \setting{Megabytes} & 3,354.28 & 355.86 & $-89.39\%$ \\
Terminal \setting{RSS_kb peak} & 3,603,800 KiB & 1,492,656 KiB & $-58.58\%$ \\
User CPU & 5,110.38 s & 15,144.10 s & $2.963\times$ \\
Wall & 5,131 s & 15,178 s & $2.958\times$ \\
\bottomrule
\end{tabular}
\end{center}

This is valuable negative evidence: representation size fell radically but CPU
nearly tripled.  The compact output reports 19.190 billion unit-conflict exact
tests, 24.503 billion backward-demodulation posting groups examined, 21.342
billion occurrences examined, and 494.44 million shallow path checks.  Both
runs made about 6.10 billion demodulation attempts.

The trajectories are not identical.  The files contain no external
\setting{/usr/bin/time} result, but their terminal \setting{Allocator_slabs}
lines contain a process peak-RSS diagnostic: 3,603,800 versus 1,492,656 KiB,
or 58.58\% lower (2.414 times smaller).  This is the appropriate process-RSS
comparison for the pair.  Because the compact run used mmap, its internal
355.86-MiB counter omits resident mapped pages and the 89.39\% value is only an
internal-accounting reduction.  Even RSS excludes unmapped filesystem cache,
so a cgroup total is still needed for whole-job memory.  Both numbers describe
an obsolete, slow build, not the performance of the accepted code.

The final implementation addresses the exposed mechanisms through shallow-path
selectivity, grouped and delta-coded occurrences, bounded 4096-ID archive
snapshots, dense stable-ID maps, coordinated term compaction, and explicit file
backing.  The exact 11,000-given input should now be rerun with the accepted
binary.

\subsection{Expected saving on very large AIM runs}

For trajectory-sensitive compact OTTER, the accepted measurement is 76.85\%
less peak RSS.  The fixed 88,494-hint bank and 131,001-clause final SOS make a
literal 80\% difficult on this particular proof.  The obsolete 11,000-given
result suggests that savings grow with passive/index dominance, but is not an
RSS basis for quantitative extrapolation.

For bounded collective DISCOUNT the intended growth law is stronger:

\begin{lstlisting}
ordinary OTTER RAM
  = fixed hints + O(all passive bodies and eager indexes) + proof history

DISCOUNT/collective RAM
  = packed hints + O(active indexes and history/descriptors)
    + O(configured candidate window) + file-backed bodies
\end{lstlisting}

Replacing 4.59 or 8.50 million resident SOS clauses by active/history state and
a bounded window makes 80--95\% whole-process reduction plausible on the
largest passive-dominated workloads.  The 90\% midpoint is a capacity-planning
hypothesis, not an accepted measurement.

The defensible statements are:
\begin{itemize}
  \item accepted: 76.85\% lower external peak RSS at the same proof boundary
        without CPU penalty;
  \item intermediate indication: at 11,000 givens, 58.58\% lower self-reported
        peak RSS and 89.39\% lower internally accounted memory, from an obsolete
        2.963-times-slower mmap build without cgroup page-cache accounting;
  \item conditional forecast: 80--95\% less total RAM on much larger
        passive-dominated AIM searches, awaiting current-binary RSS/PSS,
        cgroup page-cache, and proof/coverage experiments.
\end{itemize}

\section{What worked and what failed}

\begin{longtable}{>{\raggedright\arraybackslash}p{0.27\textwidth}
                  >{\raggedright\arraybackslash}p{0.17\textwidth}
                  >{\raggedright\arraybackslash}p{0.47\textwidth}}
\toprule
Idea & Outcome & Explanation \\
\midrule
\endfirsthead
\toprule
Idea & Outcome & Explanation \\
\midrule
\endhead
Delete rejected generated clauses & no radical gain & Most tautologies and
forward-subsumed conclusions were already transient. \\
Delete disabled clauses & useful, insufficient & Measured 4.5--30\%; proof
parents still need storage and live SOS remains. \\
Compact proof ancestors & retained & Exact proof history without live clause
graphs; necessary but not the SOS growth law. \\
Ancestor mmap & rejected as final default & Scans and reports faulted cold pages;
a 13.1-GB mapping had 10.2 GB resident. \\
Explicit file archive & retained & One reusable I/O buffer and predictable
residency. \\
Strict DISCOUNT & alternative & Best asymptotic bound, but different
demodulator availability and Osborn trajectory. \\
Collective descriptors & alternative & Bounded and checkpointable, but hint
counts did not recover the proof. \\
Eager-interreduced DISCOUNT & not default & Strong contraction, costly rewrite
debt, and changed hint normal forms. \\
First packed hints & replaced & Memory saving with severe broad-posting CPU
cost. \\
Packed-fast hints & retained & Exact trace and bounded, selective retrieval. \\
Compact bodies only & insufficient & Rewrite/unit/back indexes dominated by
1000 givens. \\
First compact indexes & correct, too slow & Billions of unit and redex
candidates in the 11k run. \\
Radix paths, shallow filters, shared terms, delta streams & retained & Smaller
and more selective without changing final logical answers. \\
Large decoded cache & rejected & Avoided decodes but did not improve CPU and
raised RSS. \\
8K final query cache & rejected & Only 596 KiB saved, with worse CPU and
fragmentation. \\
16K bounded query cache & retained & 35.08\% hits and 136 million avoided
candidates within the RSS gate. \\
Full-array compaction & replaced & Temporary copies raised peak memory. \\
Streamed/file-sorted compaction & retained & Bounds transient memory and
preserves deterministic order. \\
Release before proof replay & retained & Avoids overlapping the search state
and final proof DAG. \\
\bottomrule
\end{longtable}

The general lessons are that indexes must be judged by candidate volume rather
than bytes per node; contraction timing is part of practical algebraic search;
and resident memory is a lifetime property, not merely a sum of steady-state
record sizes.

\section{Using the accepted mode}

\subsection{Build and test}

\begin{lstlisting}[language=bash]
make all -j4
make test1
make compact-frontier-tests
make discount-tests
\end{lstlisting}

For publication record the commit, binary/input SHA-256, linked libraries, and
host.  The accepted binary SHA-256 was
\hashval{78b5ef4b2e53fc5ae2ab81cc46e6455213128e8ac083d6fb11259ad7b1339a6b};
the accepted input SHA-256 was
\hashval{9781ee07691bc62e01f67534208620ca0be3f026f55248a227e9161f1ed17e6c}.

\subsection{Compact-OTTER input}

Place the following after automatic and older experimental assignments:

\begin{lstlisting}
assign(search_loop,otter).
assign(passive_store,dense).
assign(hint_index,packed_fast).
assign(inference_frontier,clauses).
assign(ancestor_store,file).
assign(sos_limit,-1).

set(process_initial_sos).
set(back_demod).
set(back_demod_hints).
clear(unit_deletion).
clear(ancestor_subsume).
clear(eval_rewrite).
set(compact_otter_demodulation).
set(compact_otter_unit_index).
set(compact_otter_back_demod_index).
set(compact_otter_nonunit_index).

assign(compact_passive_cache,0).
assign(compact_index_stale_pct,10).
assign(compact_term_reclaim_kb,2048).
\end{lstlisting}

Start the process with the measured heap policy:

\begin{lstlisting}[language=bash]
P9_COMPACT_HEAP=1 /usr/bin/time -v bin/prover9 \
  -f your-compact-input.in \
  > your-compact-output.out \
  2> your-compact-output.time
\end{lstlisting}

The mode requires \setting{sos_limit=-1},
\setting{process_initial_sos}, a clauses frontier, all four compact indexes,
and backward demodulation.  The current authoritative compact unit/nonunit
indexes do not support \setting{unit_deletion} or
\setting{ancestor_subsume}; the compact rewrite bank is incompatible with
\setting{eval_rewrite}.

The ordinary control uses \setting{search_loop=otter},
\setting{passive_store=full}, \setting{hint_index=fpa},
\setting{inference_frontier=clauses}, and \setting{ancestor_store=off}, while
retaining identical rules, order, actions, selection, hints, and resource
limits.

\subsection{Bounded progression}

On a low-memory or slow host, begin with 100--300 givens and a deterministic
hint sample, then all hints at 300, 1000 givens, a time-bounded 2000--4000
prefix, and only then a full run.  Keep both internal and external limits.
Independent small runs may use separate real cores if aggregate RSS is known.
A hint sample is a debugging/A-B instrument, not a proof replacement.

\subsection{Archive location and measurement}

Current source hard-codes \setting{/tmp} in both archive constructors and does
not consult \setting{TMPDIR}.  It creates and immediately unlinks either a
\setting{/tmp/prover9-ancestors-XXXXXX} or
\setting{/tmp/prover9-passive-XXXXXX} file.  The open file is visible only
through \setting{/proc/PID/fd} and disappears on exit.  Ensure that the
filesystem containing \setting{/tmp} has enough space.

Logical size, allocated blocks, and residency can be inspected with
\setting{readlink}, \setting{stat -L}, \setting{du -hL}, and
\setting{/proc/PID/smaps}.  Changing \setting{TMPDIR} currently has no effect.

When claiming total job RAM rather than process RSS, also record cgroup-v2
\setting{memory.peak} and the file-cache fields in \setting{memory.stat}.

For each comparison report proof/stopping reason; all main clause populations;
proof length and \setting{prooftrans} result; user/system/wall time; external
maximum RSS; RSS/PSS at common given boundaries; archive logical bytes
separately; and compact-index bytes, exact tests, rebuilds, and validation
failures.  Do not use only the last active-hint \setting{matched=} snapshot or
the cumulative \setting{Generated} count.

\section{Future work}

\subsection{Current-binary large-profile replay}

The first priority is a replay of the exact 18,306-hint input behind the new
11,000-given pair, using file backing, compact heap policy, and external RSS/PSS.
If possible compare ordinary indexes/full bodies, final compact indexes/full
bodies, and the complete file-backed mode.  This separates index CPU from
archive I/O.

\subsection{Multi-million-SOS validation}

Run current compact OTTER and DISCOUNT/collective on several passive-dominated
problems including Osborn and AAPERM, with identical hashes and limits.  Report
equal-given prefixes and equal-resource solved coverage.  A radical claim for
these workloads should require at least 80\% lower external peak RSS,
corresponding cgroup whole-job/page-cache measurements, and at least 95\% PSS
accounting.

\subsection{Reduce exact-test amplification}

Future indexes should be driven by per-query candidate distributions.  Options
include measured selective path features, substitution/code-tree retrieval,
block-compressed posting intersections with skips, operation-specific unit
structures, and direct exact matching over serialized terms.  Every prototype
must pass the event oracle and be compared with the final index, not the first
packed version.

\subsection{True cross-clause term sharing}

The measured duplicate-subterm rate justifies a canonical immutable term DAG
prototype.  It needs symbol/child-ID keys, compact clause roots, a bounded
interning table, liveness compatible with index compaction, direct matching and
rewriting over term IDs, portable archive encoding, and a peak analysis that
includes both generations.  The current instrumentation is an upper bound on
savings, not proof that hash-consing wins.

\subsection{File-backed cold posting tiers}

Compact OTTER still retains candidate indexes for every passive in RAM.  A
two-tier external-memory index could use an in-memory directory/hot tail and
immutable compressed posting blocks on disk, with skip metadata and sequential
merge/compaction.  Backward demodulation is the difficult access pattern;
batching new rules and scanning compressed blocks may outperform millions of
small reads.

\subsection{Direct compact rewriting}

A compact rewrite pipeline could decode into a bounded workspace, rewrite, and
emit a new archived record without constructing a full \setting{Topform}
forest.  It must preserve justifications, attributes, hints, orientation, and
action side effects, and should target the backward-demodulation/preprocessing
clocks rather than only allocation counts.

\subsection{Collective proof guidance}

Collective search should use proof-relevant ancestor progress, not total hint
matches.  Candidate signals include proof-hint parent closure, distance to
unmatched successor hints, per-rule/direction quotas from source proofs, and
normal-form agreement.  A heuristic lane may reorder finite work, but fair FIFO
service and exact committed processing must remain.

\subsection{Two-stage workflow and operational hardening}

Where exact trajectory is not required, many low-memory DISCOUNT/collective
jobs can discover proofs or hints, followed by a smaller number of compact-
OTTER proof-producing replays.  Evaluate solved problems per machine-day and
verified proofs per GiB.

Operational work should honor \setting{TMPDIR} or add an archive-directory
option, check free space, expose diagnostic FD information, collect cgroup-v2
counters, record hashes and external time automatically, test the tight RSS
gate across platforms, and publish option incompatibilities.

\section{Reproducibility map}

The central source files are \setting{provers.src/search.c},
\setting{provers.src/giv_select.c},
\setting{provers.src/cold_passive_store.c},
\setting{ladr/clause_store.c}, \setting{ladr/hints.c},
\setting{ladr/hint_postings.c}, and the
\setting{compact_rewrite}, \setting{compact_unit_index},
\setting{compact_back_demod}, \setting{compact_feature_index},
\setting{compact_term_pool}, and \setting{compact_id_map} modules.

The authoritative records are
\setting{P9-PHASE5-ACCEPTANCE-AUDIT.md},
\setting{P9-PHASE5-COMPACT-FRONTIER-PLAN.md},
\setting{P9-CHAT-TEST-REPORT.md},
\setting{P9-RADICAL-RAM-REPORT.md},
\setting{P9-DISCOUNT-WALDMEISTER-PLAN.md},
\setting{P9-COLLECTIVE-SCHEDULER-PLAN.md},
\setting{P9-COLLECTIVE-DEMODULATION-PLAN.md},
\setting{P9-BETTER-PACKED-PLAN.md}, and
\setting{Checkpoint-Format-Spec.txt}.

Full-proof artifacts are under
\setting{/project/phase5-results/phase5-completion-audit/final-proof-current}.
The work from the frozen baseline to the final audit comprises 192 small
commits with detailed messages covering ancestor storage, allocation, DISCOUNT,
collective continuations, hints, compact indexes, term sharing, file backing,
bounded rebuilds, proof lifetime, and checkpoint audit.

\section{Conclusion}

Millions of SOS clauses were not merely proof debris.  Under the OTTER loop
they remained active for contraction and therefore occurred in several term
and clause indexes.  Hints and proof ancestors added fixed and historical
costs.  Radical memory reduction required changing ownership across the whole
saturation state rather than deleting one list.

The DISCOUNT/collective architecture gives the cleanest asymptotic answer and
remains promising for large passive-dominated exploration.  It also showed that
the practical behavior of an ATP is not determined by inference rules alone:
the timing of demodulation, normalization, hint matching, and rule interleaving
can decide whether a known proof is found.

Compact OTTER resolves that tension for the current trajectory-sensitive
problem.  It keeps eager contraction and hint semantics while replacing
pointer identity by stable clause numbers, resident passive bodies by an exact
file archive, and pointer-heavy indexes by compact authoritative filters.  It
proves at the same 2945-given boundary with 76.85\% less peak RSS and no CPU
penalty.

The literal 80--90\% target is not yet a universal measured claim.  The
accepted proof saves 76.85\%; the obsolete slow mmap build saved 58.58\% by its
terminal peak-RSS diagnostic while its allocator-only counter fell 89.39\%.
Current-binary 11,000-given and multi-million-SOS runs with cgroup accounting
are the next decisive experiments.

\begin{thebibliography}{12}

\bibitem{bachmair-ganzinger}
L. Bachmair and H. Ganzinger,
``Resolution Theorem Proving,'' in J. A. Robinson and A. Voronkov (eds.),
\emph{Handbook of Automated Reasoning}, vol. 1, Elsevier/MIT Press, 2001.

\bibitem{harrison}
J. Harrison, \emph{Handbook of Practical Logic and Automated Reasoning},
Cambridge University Press, 2009.
DOI: \href{https://doi.org/10.1017/CBO9780511576430}
{10.1017/CBO9780511576430}.

\bibitem{waldmeister}
T. Hillenbrand,
``Citius altius fortius: Lessons learned from the Theorem Prover Waldmeister,''
\emph{Electronic Notes in Theoretical Computer Science}, 2003/2004.
DOI: \href{https://doi.org/10.1016/S1571-0661(04)80649-2}
{10.1016/S1571-0661(04)80649-2}.

\bibitem{next-waldmeister}
T. Hillenbrand and B. L\"ochner, ``The Next Waldmeister Loop,'' in
A. Voronkov (ed.), \emph{Proceedings of CADE-18}, pp. 486--500, Springer, 2002.

\bibitem{otter}
W. McCune, \emph{OTTER 3.3 Reference Manual}, Technical Memorandum
ANL/MCS-TM-263, Argonne National Laboratory, 2003.

\bibitem{prover9}
W. McCune, ``Prover9 and Mace4,'' Prover9 distribution and manual, 2005--2009.

\bibitem{nieuwenhuis-rubio}
R. Nieuwenhuis and A. Rubio, ``Paramodulation-Based Theorem Proving,'' in
J. A. Robinson and A. Voronkov (eds.), \emph{Handbook of Automated Reasoning},
vol. 1, Elsevier/MIT Press, 2001.

\bibitem{lrs}
A. Riazanov and A. Voronkov, ``Limited Resource Strategy in Resolution Theorem
Proving,'' \emph{Journal of Symbolic Computation} 36(1--2):101--115, 2003.

\bibitem{e}
S. Schulz, ``E---A Brainiac Theorem Prover,'' \emph{AI Communications}
15(2--3):111--126, 2002.

\bibitem{hints}
R. Veroff, ``Using Hints to Increase the Effectiveness of an Automated
Reasoning Program: Case Studies,'' \emph{Journal of Automated Reasoning}
16(3):223--239, 1996.

\bibitem{sketches}
R. Veroff, ``Solving Open Questions and Other Challenge Problems Using Proof
Sketches,'' \emph{Journal of Automated Reasoning} 27(2):157--174, 2001.

\bibitem{resonance}
L. Wos, ``The Resonance Strategy,'' \emph{Computers \& Mathematics with
Applications} 29(2):133--178, 1995.

\end{thebibliography}

\end{document}
